\documentclass[tikz,border=10pt]{standalone}
\usepackage{helvet}
\renewcommand{\familydefault}{\sfdefault}
\usepackage{tikz}
\usepackage{tikz-3dplot}
\usetikzlibrary{calc,positioning,shapes.geometric,arrows.meta,intersections,patterns,decorations.pathmorphing,backgrounds}

% Colors
\definecolor{emeraldpine}{HTML}{1A4D3E}
\definecolor{harvardcrimson}{HTML}{A51C30}
\definecolor{ethdarkblue}{HTML}{1F407A}
\definecolor{ethpetrol}{HTML}{007A87}
\definecolor{ethbronze}{HTML}{B87333}
\definecolor{ethpurple}{HTML}{5B4B8A}
\definecolor{ink}{HTML}{1A202C}
\definecolor{softink}{HTML}{4A5568}
\definecolor{paleink}{HTML}{A0AEC0}
\definecolor{cardbg}{HTML}{F8FAFC}
\definecolor{shadecolor}{HTML}{F4F6F9}

\begin{document}

\tdplotsetmaincoords{70}{110} % elevated perspective

\begin{tikzpicture}[tdplot_main_coords, font=\sffamily, >={Stealth[scale=1.2]}, scale=1.5, line join=round]
    % Define scales and dimensions
    \def\carL{4} % length (x-axis)
    \def\carW{2} % width (y-axis)
    \def\carH{1.5} % height of body
    \def\chassisH{0.6} % height of chassis from ground
    \def\totalH{2.1} % total height (\chassisH + \carH)
    
    % Draw the ground plane grid
    \draw[step=1, paleink!30, thin] (-3,-2,0) grid (7,6,0);
    
    % Obstacle (human figure / icon lying down) in blind spot
    % Placed at z=0.05
    \begin{scope}[shift={(1.5, 1, 0.05)}]
        \fill[harvardcrimson] (0,0,0) circle (0.15); % head
        \fill[harvardcrimson] (-0.3,-0.6,0) -- (0.3,-0.6,0) -- (0.3,-0.2,0) -- (-0.3,-0.2,0) -- cycle; % torso
        \fill[harvardcrimson] (-0.2,-1.2,0) -- (-0.05,-1.2,0) -- (-0.05,-0.6,0) -- (-0.2,-0.6,0) -- cycle; % leg
        \fill[harvardcrimson] (0.2,-1.2,0) -- (0.05,-1.2,0) -- (0.05,-0.6,0) -- (0.2,-0.6,0) -- cycle; % leg
    \end{scope}

    % 1. Blind spot volume (red)
    % We draw the faces that are in the back and bottom first.
    % Base:
    \filldraw[fill=harvardcrimson!30, draw=harvardcrimson, thick, opacity=0.7] (0,0,0) -- (\carL,0,0) -- (\carL,\carW,0) -- (0,\carW,0) -- cycle;
    % Face x=0:
    \fill[harvardcrimson!40, opacity=0.6] (0,0,0) -- (0,\carW,0) -- (0,\carW,\chassisH) -- (0,0,\chassisH) -- cycle;
    % Face y=2 (\carW):
    \fill[harvardcrimson!40, opacity=0.6] (0,\carW,0) -- (\carL,\carW,0) -- (\carL,\carW,\chassisH) -- (0,\carW,\chassisH) -- cycle;
    % Face x=4 (\carL):
    \fill[harvardcrimson!40, opacity=0.6] (\carL,0,0) -- (\carL,\carW,0) -- (\carL,\carW,\chassisH) -- (\carL,0,\chassisH) -- cycle;
    % Top of blind spot:
    \filldraw[fill=harvardcrimson!50, draw=harvardcrimson, thick, opacity=0.8] (0,0,\chassisH) -- (\carL,0,\chassisH) -- (\carL,\carW,\chassisH) -- (0,\carW,\chassisH) -- cycle;
    
    % Wheels
    \foreach \x/\y in {0.8/0, 3.2/0, 0.8/\carW, 3.2/\carW} {
        \begin{scope}[shift={(\x,\y,\chassisH/2)}]
            % Draw wheel as a box to give it thickness
            \filldraw[fill=ink, draw=black] (-0.3, -0.1, -\chassisH/2) -- (0.3, -0.1, -\chassisH/2) -- (0.3, -0.1, \chassisH/2) -- (-0.3, -0.1, \chassisH/2) -- cycle;
            \filldraw[fill=ink!80, draw=black] (0.3, -0.1, -\chassisH/2) -- (0.3, 0.1, -\chassisH/2) -- (0.3, 0.1, \chassisH/2) -- (0.3, -0.1, \chassisH/2) -- cycle;
            \filldraw[fill=ink, draw=black] (-0.3, 0.1, -\chassisH/2) -- (0.3, 0.1, -\chassisH/2) -- (0.3, 0.1, \chassisH/2) -- (-0.3, 0.1, \chassisH/2) -- cycle;
            \filldraw[fill=ink!90, draw=black] (-0.3, -0.1, \chassisH/2) -- (0.3, -0.1, \chassisH/2) -- (0.3, 0.1, \chassisH/2) -- (-0.3, 0.1, \chassisH/2) -- cycle;
        \end{scope}
    }
    
    % 2. Vehicle chassis & body
    % Given phi=110, visible faces are x=0, y=\carW, and z=\totalH
    % Face x=0:
    \filldraw[fill=white!80!ethdarkblue, draw=ethdarkblue, thick, opacity=0.75] (0,0,\chassisH) -- (0,\carW,\chassisH) -- (0,\carW,\totalH) -- (0,0,\totalH) -- cycle;
    
    % Face y=\carW:
    \filldraw[fill=white!70!ethdarkblue, draw=ethdarkblue, thick, opacity=0.75] (0,\carW,\chassisH) -- (\carL,\carW,\chassisH) -- (\carL,\carW,\totalH) -- (0,\carW,\totalH) -- cycle;
    
    % Face x=\carL (partially hidden by y=\carW from this angle, but we might see it slightly)
    \filldraw[fill=white!60!ethdarkblue, draw=ethdarkblue, thick, opacity=0.75] (\carL,0,\chassisH) -- (\carL,\carW,\chassisH) -- (\carL,\carW,\totalH) -- (\carL,0,\totalH) -- cycle;
    
    % Top
    \filldraw[fill=white!95!ethdarkblue, draw=ethdarkblue, thick, opacity=0.9] (0,0,\totalH) -- (\carL,0,\totalH) -- (\carL,\carW,\totalH) -- (0,\carW,\totalH) -- cycle;
    
    % Roof LiDAR
    \filldraw[fill=ethpetrol!50, draw=ethdarkblue, thick] (\carL/2-0.2, \carW/2-0.2, \totalH) -- (\carL/2+0.2, \carW/2-0.2, \totalH) -- (\carL/2+0.2, \carW/2+0.2, \totalH) -- (\carL/2-0.2, \carW/2+0.2, \totalH) -- cycle;
    % sides of lidar
    \filldraw[fill=ethpetrol!40, draw=ethdarkblue, thick] (\carL/2-0.2, \carW/2+0.2, \totalH) -- (\carL/2+0.2, \carW/2+0.2, \totalH) -- (\carL/2+0.2, \carW/2+0.2, \totalH+0.3) -- (\carL/2-0.2, \carW/2+0.2, \totalH+0.3) -- cycle;
    \filldraw[fill=ethpetrol!30, draw=ethdarkblue, thick] (\carL/2+0.2, \carW/2-0.2, \totalH) -- (\carL/2+0.2, \carW/2+0.2, \totalH) -- (\carL/2+0.2, \carW/2+0.2, \totalH+0.3) -- (\carL/2+0.2, \carW/2-0.2, \totalH+0.3) -- cycle;
    % top of lidar
    \filldraw[fill=ethpetrol, draw=ethdarkblue, thick] (\carL/2-0.2, \carW/2-0.2, \totalH+0.3) -- (\carL/2+0.2, \carW/2-0.2, \totalH+0.3) -- (\carL/2+0.2, \carW/2+0.2, \totalH+0.3) -- (\carL/2-0.2, \carW/2+0.2, \totalH+0.3) -- cycle;

    % 3. Sensor Frustums (semi-transparent, projecting outwards)
    
    % Front Camera Frustum (from front bumper)
    \draw[fill=ethpetrol!40, opacity=0.3, draw=ethpetrol, thick] (\carL, \carW/2, \totalH-0.2) -- (\carL+3, \carW/2-2.5, 0) -- (\carL+3, \carW/2+2.5, 0) -- cycle;
    \draw[fill=ethpetrol!30, opacity=0.3, draw=ethpetrol, thick] (\carL, \carW/2, \totalH-0.2) -- (\carL+3, \carW/2+2.5, 0) -- (\carL+3, \carW/2+2.5, 1) -- cycle;
    
    % Side Camera Frustum (y=\carW)
    \draw[fill=ethpetrol!40, opacity=0.3, draw=ethpetrol, thick] (\carL/2, \carW, \totalH-0.2) -- (\carL/2+2, \carW+2.5, 0) -- (\carL/2-2, \carW+2.5, 0) -- cycle;

    % Roof LiDAR Cone (projecting downwards to ground, covering a 360 ring, we draw it as 4 transparent pyramids)
    \draw[fill=ethpetrol!20, opacity=0.2, draw=ethpetrol] (\carL/2, \carW/2, \totalH+0.3) -- (\carL+3, \carW+3, 0) -- (\carL+3, -1, 0) -- cycle;
    \draw[fill=ethpetrol!20, opacity=0.2, draw=ethpetrol] (\carL/2, \carW/2, \totalH+0.3) -- (\carL+3, \carW+3, 0) -- (-2, \carW+3, 0) -- cycle;
    \draw[fill=ethpetrol!20, opacity=0.2, draw=ethpetrol] (\carL/2, \carW/2, \totalH+0.3) -- (-2, \carW+3, 0) -- (-2, -1, 0) -- cycle;

    % Labels
    \node[anchor=east, text=harvardcrimson, font=\sffamily\bfseries, align=right] at (-1, -1, 1.5) (blindspot) {Unobservable State\\(Blind Spot)};
    \draw[->, harvardcrimson, thick] (blindspot) edge[bend right=15] (0.5, 0.5, 0.3);
    
    \node[anchor=south west, text=ethpetrol, font=\sffamily\bfseries, align=left] at (\carL+0.5, \carW+1.5, 2.5) (sensor) {Sensor Coverage\\(Frustums)};
    \draw[->, ethpetrol, thick] (sensor) edge[bend left=10] (\carL+0.5, \carW+1, 0.8);

\end{tikzpicture}
\end{document}
